\documentclass[11pt,a4paper]{article}
\usepackage[utf8]{inputenc}
\usepackage[T1]{fontenc}
\usepackage[margin=1in]{geometry}
\usepackage{hyperref}

\hypersetup{
  hidelinks
}

\title{\textbf{PC accessibility crisis}}
\author{Manuel Pollazzon}
\date{}

\begin{document}
\maketitle

\begin{abstract}
Since late 2025, the global technology sector has faced an acute memory crisis
characterized by explosive price inflation in Dynamic Random Access Memory
(DRAM). Driven by enterprise demand for High-Bandwidth Memory (HBM) and server
DRAM in Artificial Intelligence (AI) data centers, major manufacturers have
reallocated manufacturing capacity away from consumer hardware. This paper
synthesizes market statistics from Gartner, GSMA, and ITU to analyze the
structural mechanics of component cost inflation, its impact on entry-level PC
and smartphone markets, and the resulting widening of the global digital
divide.
\end{abstract}

\section{Introduction and Component Inflation}
The semiconductor industry experienced a fundamental shift beginning in
October 2025. Unprecedented capital expenditure in AI infrastructure created an
insatiable demand for HBM and enterprise-grade server DRAM. A single AI server
requires between 8 and 10 times the memory capacity of a standard enterprise
server, incentivizing memory fabricators to prioritize high-margin enterprise
orders over consumer-grade DDR5 components.

Contract pricing for standard DDR5 16Gb chips surged from \$8.70 in October
2025 to \$48.00 in September 2026---representing a price surge of over 450\%
(with spot market spikes reaching up to 600\%). According to market research
from Gartner, PC memory expenses are projected to reach a peak of 23\% of the
total bill-of-materials (BOM) in 2026, up from 16\% in 2025. Combined DRAM and
SSD cost increases are expected to reach 130\% by the end of 2026.

\section{Market Contraction and Segment Erasure}
The rapid rise in component costs has severely eroded margins for budget
hardware manufacturers. Vendors are increasingly unable to absorb component
inflation, leading directly to the contraction of consumer shipments and the
extinction of low-cost hardware tiers.

Key global market impact and forecast metrics (2026--2028):
\begin{itemize}
  \item \textbf{DDR5 16Gb contract price increase:} +451\% (\$8.70 to \$48.00)
  \item \textbf{Peak memory share of PC BOM (2026):} 23\% (up from 16\% in
        2025)
  \item \textbf{Global PC shipment change (2026 vs 2025):} $-10.4\%$
  \item \textbf{Global smartphone shipment change (2026 vs 2025):} $-8.4\%$
  \item \textbf{Average PC retail price increase (2026):} +17\%
  \item \textbf{Average smartphone retail price increase (2026):} +13\%
  \item \textbf{Extinction horizon for sub-\$500 entry PC market:} Year 2028
        (Gartner forecast)
\end{itemize}

Gartner explicitly predicts that the \textbf{sub-\$500 entry-level PC market
will disappear by 2028}, as manufacturers pivot exclusively toward premium and
enterprise offerings. Consequently, worldwide PC shipments are expected to
shrink by 10.4\% and smartphone shipments by 8.4\% in 2026. To adapt to higher
replacement costs, business users are extending PC lifespans by an average of
15\%, while consumers are delaying upgrades by 20\%.

\section{Socio-Economic Consequences and Digital Inclusion}
Hardware inflation poses a critical challenge to global digital inclusion
initiatives, particularly in developing economies where entry-level hardware
serves as the primary gateway to online services, education, and financial
systems.
\begin{itemize}
  \item \textbf{Sub-Saharan Africa Handset Burden:} GSMA reporting indicates
        that an entry-level smartphone currently absorbs \textbf{76\% of the
        average monthly income} for the bottom 40\% of the population in
        Sub-Saharan Africa. Sub-\$100 smartphone shipments in the region
        dropped by 34\% in Q2 2026.
  \item \textbf{Global Connectivity Divide:} ITU \textit{Facts and Figures
        2026} data shows that \textbf{2 billion people remain completely
        offline}. While internet usage reaches 93\% in high-income nations, it
        remains at just 32\% in low-income countries.
  \item \textbf{Financial Barriers:} In Least Developed Countries (LDCs),
        basic mobile broadband connectivity represents a median of 6.1\% of
        Gross National Income (GNI) per capita, far exceeding the UN Broadband
        Commission's 2\% target.
\end{itemize}

\section{Conclusion}
The structural reallocation of DRAM wafer capacity toward AI data centers has
created widespread hardware inflation across consumer markets. Without policy
intervention or targeted subsidies for basic educational hardware, rising
component costs will continue to push entry-level computing out of reach for
millions, threatening to exacerbate global economic inequality and broaden the
digital divide.

\section{Source}
\begin{itemize}
\item https://www.gartner.com/en/newsroom/press-releases/2026-02-26-gartner-says-surging-memory-costs-will-reduce-global-pc-and-smartphone-shipments-in-2026
\item https://www.itu.int/en/mediacentre/Pages/PR-2026-11-03-Facts-and-Figures.aspx
\item https://leadership.ng/smartphones-cost-76-of-monthly-income-in-sub-saharan-africa/\#1
\item https://www.techspot.com/news/114436-teamgroup-warns-memory-shortage-has-not-reached-peak.html
\end{itemize}

\end{document}